\documentclass[11pt]{article}

\usepackage[margin=1in]{geometry}
\usepackage{amsmath,amssymb}
\usepackage{booktabs}
\usepackage{graphicx}
\usepackage{multirow}
\usepackage{hyperref}
\usepackage[numbers,sort&compress]{natbib}
\usepackage{xcolor}
\usepackage{caption}
\usepackage{subcaption}
\usepackage{authblk}

\usepackage{placeins} % <--- ADD THIS HERE!
\usepackage{float}    % <--- ADD THIS HERE TOO!

\newcommand{\todo}[1]{\textcolor{red}{\textbf{[TODO: #1]}}}

% Render a figure if its file is present, otherwise a labelled placeholder so the
% document always compiles. The EDA/visualization PNGs are produced by the study
% notebook and saved to results/figures/ on Drive; copy them into paper/figures/
% (see paper/README) to replace the placeholders with the real plots.
\newcommand{\figorbox}[2][width=\linewidth]{%
  \IfFileExists{#2}%
    {\includegraphics[#1]{#2}}%
    {\framebox[0.9\linewidth]{\parbox{0.82\linewidth}{\centering\vspace{1.6em}%
      \ttfamily\small\detokenize{#2}\\[0.5em]%
      \normalfont\itshape figure file not committed; copy from Drive \texttt{results/figures/} (see paper/README)%
      \vspace{1.6em}}}}%
}

\title{\textbf{Benchmarking Canonical Deep Forecasters on BuildingsBench:\\
The Value of Weather Inputs for Short-Term Load Forecasting}}

\author[1]{\todo{Author One}}
\author[1]{\todo{Author Two}}
\affil[1]{\todo{Affiliation}}
\date{\today}
\begin{document}
\maketitle

\begin{abstract}
Short-term load forecasting (STLF) underpins reliable and economical power
system operation, and the recent BuildingsBench benchmark has standardized its
evaluation over a large-scale corpus of 900{,}000 simulated buildings together
with seven real-world datasets. However, its public leaderboard fixes a
particular model family, while the time-series forecasting literature has since
converged on a different set of standard architectures that have never been
evaluated under the same protocol. In this work we re-implement BuildingsBench's
exact preprocessing, temporal validation split, and median-aggregated metrics,
and benchmark a broad set of canonical published forecasters---PatchTST,
iTransformer, TimeXer, Informer, Crossformer, and DLinear, alongside recurrent
(LSTM, GRU) and gradient-boosted (XGBoost, LightGBM) baselines---each transformer
sized to a common $\geq$7M-parameter floor so that model capacity is never a
confound. On the official simulated test set, every trained model significantly
outperforms all three persistence baselines (Friedman $p\approx0$), and our
pipeline reproduces the paper's own persistence values to within $0.4$ NRMSE
points, validating the re-implementation. On the seven real-world datasets
(1{,}923 buildings, zero-shot), our strongest model (TimeXer, $12.57$\%
commercial NRMSE) surpasses the value reported for the original benchmark's
largest pretrained transformer. Finally, contrasting no-weather and +weather
conditions on the simulated set, we quantify \emph{how much} weather inputs help
each architecture, finding the largest gains for cross-dimension attention models
(Crossformer: $-21\%$ commercial NRMSE) and little change for recurrent, linear,
and inverted-attention models. To test whether this value is genuine or an
artifact of the physics simulator, we build a validated full-seven-channel
real-weather track---fetching ERA5 reanalysis for each real building's true
location and gating it against the benchmark's own published series before
use---and find that the simulated weather gains almost entirely vanish on real
meters (mean absolute commercial-NRMSE change $\approx0.1$ point, mixed sign).
This indicates that the exploitable weather signal is largely a property of the
simulated pretraining corpus rather than of real building load, a caution for
weather-conditioned benchmarking on synthetic data.
\end{abstract}

\section{Introduction}
\label{sec:intro}
Accurate short-term load forecasting (STLF) is central to power-system
scheduling, demand response, and the integration of variable renewable
generation. As demand-side behavior diversifies, forecasting must exploit not
only historical load but also exogenous drivers such as weather and calendar
effects \citep{buildingsbench2023}. Deep learning has become the dominant
approach, but progress has historically been fragmented across proprietary
datasets and inconsistent evaluation protocols, making cross-study comparison
difficult.

BuildingsBench \citep{buildingsbench2023} addressed this by releasing
Buildings-900K---a corpus of roughly 900{,}000 simulated buildings for
pretraining---together with a standardized zero-shot evaluation over seven
real-world datasets, fixed preprocessing, and a published leaderboard. Its
headline models are a purpose-built Transformer with a Gaussian or tokenized
output head. Since its release, however, the general time-series forecasting
community has converged on a different set of standard architectures (PatchTST,
iTransformer, TimeXer, and others) that have not been evaluated under the
BuildingsBench protocol. Two questions therefore remain open: (i) how do these
now-canonical forecasters perform on this benchmark relative to persistence and
to the benchmark's own model family, and (ii) how much do the benchmark's
weather inputs actually contribute for each architecture?

This paper makes the following contributions:
\begin{enumerate}
  \item \textbf{A faithful re-implementation} of BuildingsBench preprocessing,
  temporal train/validation split, and median-aggregated NRMSE/RPS metrics,
  verified by reproducing the paper's own persistence baselines to within
  $0.4$ NRMSE points (Section~\ref{sec:benchmark}).
  \item \textbf{A broad benchmark} of canonical deep forecasters (PatchTST,
  iTransformer, TimeXer, Informer, Crossformer, DLinear) and classical
  baselines (LSTM, GRU, XGBoost, LightGBM), each transformer sized to a common
  $\geq7\times10^{6}$-parameter budget, evaluated on both the simulated test set
  and the seven real-world datasets.
  \item \textbf{A quantitative study of weather value, and a simulation-to-real
  transfer test.} Contrasting no-weather and +weather conditions on the simulated
  test set, we measure the improvement attributable to weather inputs per
  architecture. We then ask whether that improvement is real by introducing a
  \emph{validated full-seven-channel real-weather track}: we recover each real
  building's true location and time zone from the benchmark's own scripts,
  reconstruct the full seven-channel weather record from ERA5 reanalysis
  (recovering the five channels the benchmark never published for real
  buildings), and gate every location against the benchmark's published
  temperature series before trusting it (Section~\ref{sec:weather-real-data}).
  The large simulated weather gains do not transfer, exposing a
  synthetic-versus-real gap in weather signal (Section~\ref{sec:weather-real}).
\end{enumerate}

\section{Related Work}
\label{sec:related}

\subsection{The BuildingsBench Benchmark}
\label{sec:related-bb}
BuildingsBench \citep{buildingsbench2023} comprises Buildings-900K, a
pretraining corpus of simulated hourly load generated from the U.S.\ Department
of Energy's ResStock and ComStock building-stock models, and a zero-shot
evaluation suite of seven real building datasets. The benchmark fixes a
168-hour context and 24-hour forecast horizon, a global Box--Cox load transform,
per-channel weather standardization, and reports NRMSE (coefficient of variation
of RMSE) and a ranked probability score (RPS), aggregated as the median across
per-building values with bootstrapped confidence intervals. Its published models
are Transformer-\{S,M,L\} variants with either a Gaussian or a tokenized output
head. We adopt its preprocessing and metrics verbatim (Section~\ref{sec:method})
so that our numbers are directly comparable to its leaderboard.

\subsection{Deep Forecasting Architectures}
\label{sec:related-arch}
We benchmark architectures spanning the main design families of modern
time-series forecasting. \textbf{Patch-based transformers} such as PatchTST
\citep{patchtst} segment a series into subseries patches processed
channel-independently. \textbf{Inverted attention} (iTransformer,
\citealp{itransformer}) treats each variate as a token and attends across
variables. \textbf{Exogenous-aware transformers} such as TimeXer
\citep{timexer} explicitly separate endogenous patch tokens from exogenous
variate tokens, making them a natural fit for weather-conditioned load
forecasting. \textbf{Efficiency-oriented} designs include Informer's ProbSparse
attention \citep{informer} and Crossformer's two-stage cross-dimension attention
\citep{crossformer}. As a linear control, we include DLinear
\citep{dlinear}, which questioned whether transformers are necessary at all. We
also retain recurrent (LSTM, GRU) and gradient-boosted tree (XGBoost, LightGBM)
baselines, which remain competitive in operational STLF. Each transformer is
scaled to a common $\geq7$M-parameter budget so that architectural differences,
rather than capacity, drive the comparison.

\section{Methodology}
\label{sec:method}

\subsection{Dataset and Test Sets}
\label{sec:dataset}
We use a compute-bounded subsample of $N=20{,}000$ buildings from Buildings-900K
for training, drawn via the official weekly-sampled index file, and evaluate on
two held-out sets: (i) the official \emph{Buildings-900K-test} simulated split
($N=2{,}928$ buildings), and (ii) a real-world zero-shot set of $1{,}923$
buildings ($970$ commercial, $953$ residential) pooled from seven open datasets
(BDG-2, Borealis, Electricity, IDEAL, LCL, SMART, Sceaux); the $1{,}922$ of these
with at least one full $168{+}24$\,h window enter the metrics. Multi-year files
of the same physical building are concatenated into a single continuous series,
matching the official loader; our building count agrees with the benchmark's
reported ``over 1{,}900'' real buildings.

Figure~\ref{fig:eda} characterizes the data. The load heatmaps
(Figure~\ref{fig:eda-heatmap}) show the strong diurnal and seasonal structure
that forecasting must capture, and the differing commercial/residential
profiles; Figure~\ref{fig:eda-corr} shows the per-channel weather--load
correlation, motivating the weather-value analysis of
Section~\ref{sec:weather}.

\begin{figure}[t]
  \centering
  \begin{subfigure}[b]{0.9\textwidth}
    \figorbox[width=\textwidth]{figures/load_heatmap.png}
    \caption{Mean normalized load by day-of-year (x) and hour-of-day (y),
    commercial vs.\ residential.}
    \label{fig:eda-heatmap}
  \end{subfigure}
  \\[0.5em]
  \begin{subfigure}[b]{0.48\textwidth}
    \figorbox[width=\textwidth]{figures/daily_profiles.png}
    \caption{Median daily profile with IQR band.}
    \label{fig:eda-profile}
  \end{subfigure}
  \hfill
  \begin{subfigure}[b]{0.48\textwidth}
    \figorbox[width=\textwidth]{figures/weather_load_corr.png}
    \caption{Weather--load correlation per channel.}
    \label{fig:eda-corr}
  \end{subfigure}
  \caption{Exploratory analysis of the training subsample and real test set.
  }
  \label{fig:eda}
\end{figure}

\subsection{Data Preprocessing}
\label{sec:preprocessing}
We use the benchmark's own already-fit transforms rather than re-fitting them: a
global Box--Cox transform on load and a per-channel standardization of the seven
official weather channels (temperature, humidity, wind speed, wind direction,
and global horizontal, direct normal, and diffuse horizontal radiation). Both
are downloaded from the benchmark's public artifact bucket and validated at load
time by a numeric round-trip check. Calendar features (day-of-year, day-of-week,
hour-of-day) are linearly scaled to $[-1,1]$, matching the official transform
(not sinusoidally encoded). The context length is $L=168$h and the forecast
horizon $H=24$h. We note for reproducibility that Buildings-900K parquet rows
are not stored chronologically (confirmed empirically: $>99.9\%$ of adjacent
rows are non-consecutive in time); every window is therefore built from
timestamp-sorted, monotonicity-checked series.

\subsection{Real-Building Weather Acquisition and Validation}
\label{sec:weather-real-data}
Evaluating a +weather model on real buildings requires the same seven weather
channels the model was trained on, but the benchmark's public bucket publishes
only \emph{temperature} (and a derived humidity) for its real datasets---the
other five channels (wind speed, wind direction, and the three radiation
components) exist nowhere for the real buildings. Its own real-building weather
path therefore uses temperature alone. To evaluate our +weather checkpoints
under the exact input they expect, we reconstruct the full seven-channel record
for every real building from a public reanalysis archive, and---because a
weather series silently misaligned in location, time zone, or units would
corrupt the entire comparison---we admit it only after an explicit validation
gate.

\paragraph{Source and alignment.} We fetch all seven channels from Open-Meteo's
ERA5 archive \citep{openmeteo,era5}, the same reanalysis family the benchmark
authors used. Crucially, we do not guess coordinates: for each of the ten
weather locations (one per real dataset, plus one per BDG-2 campus site) we take
the latitude/longitude, date range, and fixed local-standard-time (no-DST)
offset \emph{verbatim from the benchmark's own weather-download script}, so our
series is anchored to the buildings' true locations rather than the CONUS proxy
coordinates in the public metadata. Data are requested in UTC with wind speed in
m\,s$^{-1}$ (matching the training channels) and then shifted by the recovered
integer-hour offset to local standard time, reproducing the authors'
index-shift convention. Each location is fetched once and cached; the seven
channels are then standardized with the benchmark's own already-fit per-channel
scalers (Section~\ref{sec:preprocessing}), never re-fit on the real data.

\paragraph{Validation gate.} Before any fetched series is used, we compare its
temperature against the benchmark's \emph{published} temperature for that
location over their overlapping hours, and admit the location only if all three
checks pass: (i) Pearson $r>0.95$ at zero lag; (ii) the best cross-correlation
lag over $\pm12$\,h is exactly $0$\,h, which catches time-zone errors by
detecting a shifted diurnal cycle; and (iii) the mean temperature offset is
$<3^{\circ}$C, which catches unit errors (e.g.\ Kelvin-vs-Celsius or a
wind-speed unit slip). Any location that fails is dropped from the cache and its
buildings simply do not enter the +weather track, so a misaligned series can
never silently reach the evaluation. All ten locations passed
(Table~\ref{tab:weather-gate}): temperature correlation ranged from $0.958$ to
$0.995$, every best lag was $0$\,h, and the largest mean offset was
$1.06^{\circ}$C, so all $1{,}923$ real buildings received validated
seven-channel weather.

\begin{table}[H]
\centering
\caption{Real-weather validation gate. Fetched ERA5 temperature vs.\ the
benchmark's published temperature per location: zero-lag Pearson $r$, best
cross-correlation lag over $\pm12$\,h (a nonzero lag would flag a time-zone
error), and mean absolute temperature offset (a large value would flag a unit
error). All ten locations pass ($r>0.95$, lag $=0$, offset $<3^{\circ}$C).}
\label{tab:weather-gate}
\small
\begin{tabular}{llrrr}
\toprule
Location (dataset/site) & Region & $r$@lag\,0 & Best lag (h) & $|\Delta T|$ ($^\circ$C) \\
\midrule
LCL         & London, UK        & 0.994 & 0 & 0.05 \\
IDEAL       & Edinburgh, UK     & 0.986 & 0 & 0.11 \\
Sceaux      & Sceaux, FR        & 0.994 & 0 & 0.17 \\
Borealis    & Waterloo, CA      & 0.993 & 0 & 0.17 \\
Electricity & Lisbon, PT        & 0.991 & 0 & 0.01 \\
SMART       & Amherst, US       & 0.995 & 0 & 0.27 \\
BDG-2 / Rat     & Washington DC, US & 0.995 & 0 & 0.47 \\
BDG-2 / Fox     & Tempe, US         & 0.991 & 0 & 0.08 \\
BDG-2 / Panther & Orlando, US       & 0.958 & 0 & 0.98 \\
BDG-2 / Bear    & Berkeley, US      & 0.958 & 0 & 1.06 \\
\bottomrule
\end{tabular}
\end{table}

\subsection{Training}
\label{sec:training}
Each neural model maps the RevIN-normalized, Box--Cox-transformed load history
and exogenous features to a Gaussian predictive distribution $(\mu_n,\sigma_n)$
over the 24-hour horizon, trained with the Gaussian negative log-likelihood
using AdamW and a cosine learning-rate schedule with warmup, for up to $200$
epochs with early stopping (patience $20$). Following the benchmark, validation
uses a temporal holdout of the final two weeks of the year on the AMY2018-release
buildings (a genuine chronological year), rather than a random building-level
split. Every transformer baseline is sized to $\geq7\times10^{6}$ parameters in
both conditions (Table~\ref{tab:compute}) so that weaker performance cannot be
attributed to insufficient capacity. All models are evaluated under two
conditions: \textbf{A} (no weather; load and calendar only) and \textbf{B}
(+weather; the seven official channels added as exogenous inputs).

\subsection{Evaluation}
\label{sec:evaluation}
We report the benchmark's exact metrics \citep{buildingsbench2023}. The point
metric is NRMSE (coefficient of variation of RMSE), computed per building as
\begin{equation}
\mathrm{NRMSE}_b = 100 \times
\frac{\sqrt{\frac{1}{n_b}\sum_i (\hat{y}_{b,i}-y_{b,i})^2}}
     {\frac{1}{n_b}\sum_i y_{b,i}},
\label{eq:nrmse}
\end{equation}
and aggregated as the \emph{median across buildings} within each building type
(commercial/residential), with 95\% bootstrap confidence intervals over $2{,}000$
resamples. The probabilistic metric (RPS) is the Gaussian continuous ranked
probability score, reported in physical kWh via the benchmark's inverse-Box--Cox
Gaussian approximation. Pairwise significance against the average-persistence
baseline uses a paired Wilcoxon signed-rank test with Holm--Bonferroni
correction; overall model-set differences use a Friedman test with Nemenyi
post-hoc ranking. We report both conditions on both test sets: condition A
(no weather) uses the condition-A checkpoints, and condition B (+weather) uses
the condition-B checkpoints fed the seven official channels---on the simulated
set from the benchmark's own weather, and on the real set from the validated
full seven-channel ERA5 record of Section~\ref{sec:weather-real-data}.

\section{Results}
\label{sec:results}

\subsection{Training Dynamics and Computational Cost}
\label{sec:compute}
Figure~\ref{fig:loss} shows validation-loss trajectories. All models converge
smoothly; notably, the +weather (condition B) curves reach substantially lower
validation NLL than their no-weather counterparts for attention-based models
(e.g.\ Crossformer B reaches $0.16$ vs.\ $0.67$ for A), an early indication that
weather inputs carry exploitable signal for these architectures.

\begin{figure}[t]
  \centering
  \figorbox[width=0.72\textwidth]{figures/loss_curves.png}
  \caption{Validation Gaussian NLL vs.\ epoch (solid = +weather, dashed =
  no-weather), from training logs. Attention models (Crossformer) show a much
  larger no-weather$\to$+weather gap than recurrent models (LSTM).}
  \label{fig:loss}
\end{figure}

Table~\ref{tab:compute} summarizes model size and training cost. Persistence
baselines are closed-form (no parameters, no training). Transformer baselines
range from $7.6$M (PatchTST) to $17.1$M (Crossformer/B) parameters and train in
$0.3$--$4.1$k seconds on a single NVIDIA A100-SXM4-40GB; the linear (DLinear)
and recurrent (LSTM, GRU) models are one to two orders of magnitude smaller.

\begin{table}[H]
\centering
\caption{Model size and training cost (single A100-SXM4-40GB). Persistence
baselines are closed-form. A = no weather, B = +weather.}
\label{tab:compute}
\begin{tabular}{llrrrr}
\toprule
Model & Cond.\ & Params & Epochs & Train (s) & Eval (s) \\
\midrule
LSTM         & A & 131{,}472    & 200 & 905    & 21.0 \\
LSTM         & B & 147{,}600    & 200 & 924    & 21.2 \\
GRU          & A & 103{,}344    & 160 & 842    & 21.4 \\
GRU          & B & 119{,}472    & 200 & 1{,}070 & 21.6 \\
DLinear      & A & 60{,}640     & 104 & 370    & 16.2 \\
DLinear      & B & 103{,}648    & 135 & 494    & 16.5 \\
PatchTST     & A & 7{,}557{,}936 & 119 & 1{,}302 & 26.6 \\
PatchTST     & B & 7{,}622{,}448 & 200 & 2{,}210 & 26.4 \\
iTransformer & A & 8{,}782{,}896 & 81  & 658    & 20.8 \\
iTransformer & B & 9{,}388{,}592 & 39  & 349    & 23.7 \\
TimeXer      & A & 7{,}581{,}552 & 189 & 1{,}804 & 22.5 \\
TimeXer      & B & 8{,}186{,}800 & 200 & 2{,}005 & 22.7 \\
Informer     & A & 8{,}023{,}344 & 123 & 2{,}143 & 50.4 \\
Informer     & B & 8{,}087{,}856 & 200 & 3{,}494 & 45.3 \\
Crossformer  & A & 11{,}963{,}568 & 86  & 1{,}633 & 39.4 \\
Crossformer  & B & 17{,}050{,}608 & 152 & 4{,}125 & 58.7 \\
\bottomrule
\end{tabular}
\end{table}


\subsection{Benchmark Results}
\label{sec:benchmark}
Table~\ref{tab:benchmark} reports the main benchmark. Our pipeline reproduces
the benchmark's own persistence values closely (Average Persistence: $33.51$
here vs.\ $33.10$ published on the simulated set), validating the
re-implementation. \textbf{On the simulated test set}, every trained model
significantly beats \texttt{persistence\_avg} under Holm--Bonferroni correction
in both conditions. The best commercial NRMSE is achieved by Crossformer with
weather ($14.68$), followed by TimeXer ($16.43$); without weather, PatchTST
($18.36$) and TimeXer ($18.45$) lead. \textbf{On the real-world zero-shot set}
(no-weather), the ordering shifts: TimeXer is strongest ($12.57$ commercial),
narrowly ahead of iTransformer ($12.68$), PatchTST ($12.78$), and Crossformer
($12.78$)---the clear simulated-set Crossformer lead does not fully transfer to
real buildings, consistent with the real data being out-of-distribution
relative to the simulated pretraining corpus.

\begin{table}[H]
\centering
\caption{Benchmark results: median NRMSE (\%) and RPS (Gaussian CRPS, kWh) across
buildings, by type. \textbf{A}=no weather, \textbf{B}=+weather. Simulated results
are shown for both conditions; the real column here is the no-weather headline
over every real building, and the +weather real-building results---evaluated with
the validated ERA5 weather of Section~\ref{sec:weather-real-data}---are reported
separately in Table~\ref{tab:weather-real}. Best in each column in
\textbf{bold}. RPS is undefined (--) for the copy-persistence baselines, which
produce no predictive variance.}
\label{tab:benchmark}
\small
\begin{tabular}{ll rr rr rr}
\toprule
& & \multicolumn{4}{c}{\textbf{Simulated (Buildings-900K-test)}} & \multicolumn{2}{c}{\textbf{Real (zero-shot)}} \\
\cmidrule(lr){3-6}\cmidrule(lr){7-8}
& & \multicolumn{2}{c}{Commercial} & \multicolumn{2}{c}{Residential} & \multicolumn{2}{c}{Com.\ / Res.} \\
Model & Cond.\ & NRMSE & RPS & NRMSE & RPS & NRMSE & RPS \\
\midrule
\multirow{2}{*}{Persistence (avg)}
     & A & 33.51 & 5.70 & 55.46 & 0.72 & 16.44 / 78.02 & 6.83 / 0.069 \\
     & B & 33.51 & 5.70 & 55.46 & 0.72 & -- & -- \\
\multirow{2}{*}{Persistence (last day)}
     & A & 35.06 & -- & 58.91 & -- & 16.43 / 99.08 & -- \\
     & B & 35.06 & -- & 58.91 & -- & -- & -- \\
\multirow{2}{*}{Persistence (last week)}
     & A & 32.79 & -- & 74.57 & -- & 19.32 / 100.24 & -- \\
     & B & 32.79 & -- & 74.57 & -- & -- & -- \\
\midrule
\multirow{2}{*}{LSTM}
     & A & 21.17 & 3.66 & 47.44 & 0.62 & 13.99 / 80.40 & 5.60 / 0.071 \\
     & B & 20.80 & 3.64 & 46.07 & 0.59 & -- & -- \\
\multirow{2}{*}{GRU}
     & A & 21.54 & 3.74 & 47.42 & 0.62 & 14.18 / 80.41 & 5.65 / 0.071 \\
     & B & 20.95 & 3.67 & 46.34 & 0.60 & -- & -- \\
\multirow{2}{*}{DLinear}
     & A & 25.64 & 4.08 & 48.87 & 0.67 & 13.68 / 80.92 & 5.47 / 0.076 \\
     & B & 25.26 & 4.07 & 48.06 & 0.65 & -- & -- \\
\multirow{2}{*}{XGBoost}
     & A & 22.31 & 5.23 & 47.41 & 0.64 &  14.17 / 79.63 &  \\
     & B & 21.69 & 5.08 & 46.56 & 0.63 & -- & -- \\
\multirow{2}{*}{LightGBM}
     & A & 21.00 & 4.95 & 46.88 & 0.64 & 13.56 / 79.53&  \\
     & B & 20.39 & 4.80 & 45.50 & 0.62 & -- & -- \\
\multirow{2}{*}{PatchTST}
     & A & 18.36 & 2.74 & 44.09 & 0.58 & 12.78 / 78.29 & 4.93 / 0.069 \\
     & B & 17.53 & 2.55 & 42.15 & 0.55 & -- & -- \\
\multirow{2}{*}{iTransformer}
     & A & 19.07 & 3.00 & 45.53 & 0.59 & 12.68 / 78.89 & 4.96 / 0.069 \\
     & B & 18.70 & 3.10 & 44.54 & 0.58 & -- & -- \\
\multirow{2}{*}{TimeXer}
     & A & 18.45 & 2.71 & 44.03 & 0.57 & \textbf{12.57} / \textbf{77.81} & 4.83 / 0.067 \\
     & B & 16.43 & 2.46 & 41.22 & 0.52 & -- & -- \\
\multirow{2}{*}{Informer}
     & A & 18.59 & 2.84 & 44.37 & 0.58 & 13.20 / 79.07 & 5.06 / 0.069 \\
     & B & 17.54 & 2.61 & 42.32 & 0.55 & -- & -- \\
\multirow{2}{*}{Crossformer}
     & A & 18.65 & 2.85 & 44.49 & 0.58 & 12.78 / 78.67 & 4.88 / 0.068 \\
     & B & \textbf{14.68} & \textbf{2.18} & \textbf{40.33} & \textbf{0.50} & -- & -- \\
\bottomrule
\end{tabular}
\end{table}

\textbf{Statistical significance.} A Friedman test across the full model set
rejects the null of no difference in both conditions ($p\approx0$; Nemenyi
critical difference $\mathrm{CD}=0.456$). The average ranks under condition B
(lower is better) are: Crossformer $1.58$, TimeXer $2.69$, PatchTST $2.75$,
Informer $3.24$, iTransformer $4.90$, LSTM $6.80$, GRU $7.18$, DLinear $7.31$,
and the three persistence variants worst ($9.57$--$10.18$). The attention-based
transformers thus form a clearly separated top tier.

\subsection{Comparison of Weather versus No-Weather}
\label{sec:weather}
Table~\ref{tab:weather} isolates the effect of adding weather inputs on the
simulated test set; Section~\ref{sec:weather-real} then asks whether this effect
survives on real buildings. Weather helps every trained model here, but the
magnitude varies sharply by architecture. Cross-dimension and exogenous-aware attention models benefit most
(Crossformer $-21.3\%$, TimeXer $-10.9\%$ relative commercial NRMSE reduction),
whereas iTransformer is nearly unchanged ($-1.9\%$) and the linear and tree
models improve only marginally ($1$--$3\%$). This mirrors the loss-curve gap in
Figure~\ref{fig:loss}: architectures designed to model cross-variable
interactions extract more from the weather channels.

\begin{table}[H]
\centering
\caption{Effect of weather inputs on the simulated test set: commercial NRMSE
under no-weather (A) and +weather (B), and relative reduction.}
\label{tab:weather}
\begin{tabular}{lrrr}
\toprule
Model & NRMSE (A) & NRMSE (B) & $\Delta$ (\%) \\
\midrule
Crossformer  & 18.65 & \textbf{14.68} & $-21.3$ \\
TimeXer      & 18.45 & 16.43 & $-10.9$ \\
Informer     & 18.59 & 17.54 & $-5.6$ \\
PatchTST     & 18.36 & 17.53 & $-4.5$ \\
LightGBM     & 21.00 & 20.39 & $-2.9$ \\
XGBoost      & 22.31 & 21.69 & $-2.8$ \\
GRU          & 21.54 & 20.95 & $-2.7$ \\
LSTM         & 21.17 & 20.80 & $-1.7$ \\
iTransformer & 19.07 & 18.70 & $-1.9$ \\
DLinear      & 25.64 & 25.26 & $-1.5$ \\
\bottomrule
\end{tabular}
\end{table}

We probed \emph{where} this value arises with a per-channel knockout analysis
(zeroing one weather channel at a time on condition-B checkpoints). Most models
were near-insensitive to any single channel (temperature-knockout
$\Delta\mathrm{NRMSE}\in[-0.08,+0.97]$), while Crossformer ($+40.2$) and TimeXer
($+15.7$) degraded sharply.

\subsection{Does Weather Transfer to Real Buildings?}
\label{sec:weather-real}
The weather gains above are large on the simulated set, but the simulated load is
itself \emph{generated} by a physics engine driven by those very weather
channels, so weather there is close to a causal input. Whether the gains reflect
usable structure in real load is a separate question, and one the benchmark's own
real-weather path (temperature only) cannot answer. We answer it directly by
evaluating each architecture's +weather (condition-B) checkpoint on the real
buildings with the validated full seven-channel ERA5 weather of
Section~\ref{sec:weather-real-data}, and comparing it against the same
architecture's no-weather (condition-A) checkpoint on the same buildings
(Table~\ref{tab:weather-real}). Because the ERA5 series passed a strict alignment
gate (Table~\ref{tab:weather-gate}), a null result here cannot be blamed on bad
weather data.

The simulated weather value does not transfer. Commercial NRMSE moves by at most
$0.30$ points across all eight architectures---mean absolute change $0.11$ points,
under $1\%$ relative---and the sign is mixed: PatchTST ($-0.13$) and Informer
($-0.20$) improve slightly, while iTransformer ($+0.30$) and most others are flat
or marginally worse. Residential NRMSE is likewise near-flat (mean absolute
change $\approx1.0$ point on a $\sim$79 base). The contrast with simulation is
stark for exactly the models that gained most there: Crossformer, which weather
improved by $21\%$ ($3.97$ points) on the simulated set, changes by $+0.02$
points on real buildings, and TimeXer's $11\%$ simulated gain becomes $+0.07$. As
an internal check, the weather-invariant persistence baseline reads identically
in both columns, confirming that the harness feeds weather only where intended.
We read this as evidence that the exploitable weather signal in Buildings-900K is
substantially a property of its physics-based generation---where weather is a
near-deterministic driver of load---rather than of real metered demand, in which
occupancy and behavior dominate once recent load history is known. This is a
caution for interpreting weather-conditioned results on synthetic corpora, and it
motivates the limitation and future directions in Section~\ref{sec:limitations}.

\begin{table}[H]
\centering
\caption{Simulation-to-real weather transfer: median commercial and residential
NRMSE (\%) on the real zero-shot buildings, no-weather (condition-A checkpoints)
vs.\ +weather (condition-B checkpoints fed the validated full seven-channel ERA5
record of Section~\ref{sec:weather-real-data}). $\Delta$ = (+weather) $-$
(no weather); negative means weather helps. The large simulated-set gains
(cf.\ Table~\ref{tab:weather}) do not appear on real buildings. Persistence is
weather-invariant and shown as a harness sanity check; tree models were not run
in the +weather real track.}
\label{tab:weather-real}
\small
\begin{tabular}{l rrr rrr}
\toprule
& \multicolumn{3}{c}{Commercial NRMSE} & \multicolumn{3}{c}{Residential NRMSE} \\
\cmidrule(lr){2-4}\cmidrule(lr){5-7}
Model & No wx & +ERA5 & $\Delta$ & No wx & +ERA5 & $\Delta$ \\
\midrule
Persistence (avg) & 16.44 & 16.44 & $0.00$ & 78.02 & 78.02 & $0.00$ \\
\midrule
LSTM         & 13.99 & 14.05 & $+0.06$ & 80.40 & 81.43 & $+1.03$ \\
GRU          & 14.18 & 14.18 & $+0.01$ & 80.41 & 81.79 & $+1.38$ \\
DLinear      & 13.68 & 13.78 & $+0.10$ & 80.92 & 82.32 & $+1.40$ \\
PatchTST     & 12.78 & 12.66 & $-0.13$ & 78.29 & 79.23 & $+0.94$ \\
iTransformer & 12.68 & 12.98 & $+0.30$ & 78.89 & 77.65 & $-1.25$ \\
TimeXer      & 12.57 & 12.64 & $+0.07$ & 77.81 & 78.35 & $+0.53$ \\
Informer     & 13.20 & 13.00 & $-0.20$ & 79.07 & 80.14 & $+1.07$ \\
Crossformer  & 12.78 & 12.80 & $+0.02$ & 78.67 & 78.65 & $-0.02$ \\
\bottomrule
\end{tabular}
\end{table}

\subsection{Comparison with the Original Benchmark Model}
\label{sec:vs-original}
Table~\ref{tab:vs-paper} compares our best no-weather model on the real-world
zero-shot benchmark against the values reported for the original benchmark's
pretrained Transformer family \citep{buildingsbench2023}. Our TimeXer, trained
on a $20{,}000$-building subsample without weather, attains $12.57\%$ commercial
and $77.81\%$ residential NRMSE, both lower than the reported Transformer-L
(Gaussian). We emphasize two caveats: the original models were pretrained on the
full $\sim$900K corpus (vs.\ our $20$K subsample), and small differences may
reflect implementation and evaluation-harness details; the comparison should be
read as ``competitive with, and on these two aggregates ahead of, the reported
values,'' not as a controlled head-to-head. 

\begin{table}[H]
\centering
\caption{Real-world zero-shot NRMSE (\%): our best model (no weather) vs.\ values
reported for the original benchmark's pretrained transformers
\citep{buildingsbench2023}. Original models pretrained on the full corpus; ours
on a 20K subsample (see caveats in text).}
\label{tab:vs-paper}
\begin{tabular}{llrr}
\toprule
Source & Model & Com.\ NRMSE & Res.\ NRMSE \\
\midrule
\multirow{4}{*}{\citet{buildingsbench2023}}
  & Transformer-L (Gaussian) & 13.31 & 79.34 \\
  & Transformer-M (Gaussian) & 13.28 & 92.60 \\
  & Transformer-S (Gaussian) & 13.97 & 102.30 \\
  & Persistence Ensemble     & 16.68 & 77.88 \\
\midrule
\multirow{2}{*}{This study}
  & \textbf{TimeXer}         & \textbf{12.57} & \textbf{77.81} \\
  & iTransformer             & 12.68 & 78.89 \\
\bottomrule
\end{tabular}
\end{table}

\subsection{Visualization}
\label{sec:viz}


\begin{figure}[H]
    \centering
    \figorbox[width=0.8\textwidth]{figures/actual_versus_predicted.png}
    \caption{Daily average actual versus predicted load across a 24-hour forecast horizon for the top 10 models. All top ten models successfully capture the overarching diurnal rhythm and multi-peak structure of the actual load across the 24-hour horizon, though their precision diverges significantly during rapid load transitions. Attention-based architectures such as Crossformer and TimeXer consistently track closest to the true trajectory, effectively mirroring complex turning points like the sharp morning dip at hour 7 and the afternoon peak at hour 15.}
    \label{fig:actual_vs_predicted}
\end{figure}

% This forces LaTeX to print the table and figure above BEFORE starting the Conclusion
\FloatBarrier

\section{Limitations and Future Work}
\label{sec:limitations}
We frame the study's boundaries as scope conditions on the questions it can
answer, and pair each with the experiment that would extend it.

\paragraph{Pretraining scale and variance.} Our models are trained on a
compute-bounded $20{,}000$-building subsample of Buildings-900K under a single
seed, whereas the original benchmark pretrains on the full $\sim$900K corpus. Our
architectural comparison controls for capacity (a common $\geq7$M-parameter
floor) but not for pretraining scale, so absolute NRMSE would likely improve with
the full corpus, and we report building-level bootstrap confidence intervals
rather than seed-level variance. Full-corpus, multi-seed training is the natural
way to firm up both the absolute numbers and the rankings.

\paragraph{Model set and output head.} We report the ten architectures trained
to completion under both conditions. The benchmark's own Transformer-\{S,M,L\}
(beyond their published values), Autoformer \citep{autoformer}, and four
post-BuildingsBench architectures---a TFT-style variable-selection model
\citep{tft}, xLSTM \citep{xlstm}, a spectral--mixer hybrid of TSMixer
\citep{tsmixer} and FITS \citep{fits}, and Mamba \citep{mamba}---were only
partially trained in our compute budget and are excluded to keep the protocol
uniform; completing them would
enable a controlled, same-pipeline replacement for the published-value
comparison of Table~\ref{tab:vs-paper}. All neural models use a Gaussian
negative-log-likelihood head; the benchmark's tokenized/quantile head and
sharper calibration diagnostics beyond CRPS are left to future work.

\paragraph{The simulation-to-real weather gap.} Our clearest new finding---that
the large simulated weather gains do not transfer to real buildings
(Section~\ref{sec:weather-real})---is established with a validated weather record,
but its \emph{mechanism} is inferred rather than isolated. We attribute the null
to Buildings-900K's physics-based generation, in which weather is a
near-deterministic driver of load, whereas real demand is dominated by occupancy
and behavior once recent load is known. Testing this directly---by fine-tuning
the +weather checkpoints on real weather--load pairs, or by measuring
weather value on real extreme-temperature days where it should concentrate---is
the most important follow-up, and would establish whether real weather signal is
merely harder to exploit or genuinely weak.

\paragraph{Operational realism.} The real +weather track uses ERA5
\emph{reanalysis}, an idealized hindcast; a deployed forecaster would instead
consume numerical-weather-prediction forecasts and inherit their errors.
Substituting forecast weather at inference time, extending the horizon (our
$72$\,h autoregressive rollout widens the +weather advantage on the simulated
set and should be re-examined on real data), and adding datasets beyond the seven
studied here would together sharpen the picture of when weather inputs pay off in
practice.

\section{Conclusion}
\label{sec:conclusion}
We re-benchmarked a broad set of canonical deep forecasting architectures on
BuildingsBench under its exact preprocessing and evaluation protocol, with each
transformer sized to a common parameter budget. Three findings stand out:
\begin{enumerate}
  \item The re-implementation is faithful: closed-form persistence baselines
  reproduce the benchmark's published values to within $0.4$ NRMSE points, and
  every trained model significantly outperforms all persistence baselines on the
  simulated test set (Friedman $p\approx0$).
  \item Attention-based transformers (Crossformer, TimeXer, PatchTST) form a
  clearly separated top tier. Rankings are not identical across settings:
  Crossformer wins the simulated set with weather, while TimeXer generalizes best
  to real buildings zero-shot, where it is competitive with---and on commercial
  and residential aggregates ahead of---the values reported for the original
  benchmark's largest pretrained transformer.
  \item Weather inputs help every model on the simulated set, but unevenly:
  cross-dimension attention models gain the most ($-21\%$ commercial NRMSE for
  Crossformer), while recurrent, linear, and inverted-attention models change
  little. Crucially, this value does not survive the move to real buildings:
  evaluated with validated full seven-channel ERA5 weather, the same +weather
  checkpoints change real commercial NRMSE by under $1\%$ on average, with mixed
  sign. The exploitable weather signal is thus largely a property of the physics
  simulator's weather-driven load generation rather than of real metered demand
  ---a caution for weather-conditioned benchmarking on synthetic corpora.
\end{enumerate}


\bibliographystyle{plainnat}
\bibliography{refs}

\end{document}
